\section{Review of PubChem AID 1851}
  

PubChem AID 1851 contains data for inhibition of five major CYP isozymes (1A2, 2C9, 2C19, 2D6 and 3A4) by 17,143 chemical compounds~\cite{Veith2009}. The tested compounds were all drugs or drug-like compounds. Most of these compounds had molecular weight below 500 daltons and logP below 5~\cite{Lapins2013}.

PubChem AID 1851 used low-luminescent substrates, which are converted to more-luminescent metabolites as the signal for this assay. Reaction progress was measured as increasing light intensity during CYP metabolism of the substrate. Inhibitors of a particular CYP isozyme reduce the rate of substrate metabolism and thus decrease the luminescent signal~\cite{Zlokarnik2005}.

The most recent CYP inhibition assay technology is based on substrates that release luciferin as the metabolite. This is a coupled assay system in which addition of luciferase and ATP converts the freed luciferin to des-carboxyluciferin along with light emission. The format is similar to fluorescence methods, but particularly well-suited to high-throughput applications because all that is required are addition-only manipulations and luminescence plate readers~\cite{Zlokarnik2005}. The assay obtained luminescence readings at a range of compound concentrations and then determined activity parameters using the Hill equation.

In PubChem AID 1851, compounds are classified as inhibitors or non-inhibitors for each CYP with an activity cutoff set to AC50 = 10~$\mu$M (AC50, “activity concentration 50”, refers to the concentration that is required to elicit half-maximal effect). However, in cases where the dose-response curve for a compound showed poor fit or the inhibition efficacy was below 60\%, the assay results were regarded as inconclusive~\cite{Veith2009}.

Each compound received an Activity Score from 0 to 100, which combines its AC50 with confidence measures for dose-response curve completeness and inhibition efficacy. Higher scores indicate stronger inhibition, higher confidence, or both. Compounds scoring 40--100 are inhibitors, those scoring 0 are non-inhibitors, and those scoring 1--39 are inconclusive~\cite{Lapins2013}. In this study, inconclusive compounds were not discarded: every compound with an Activity Score below 40 was labelled a non-inhibitor for binary classification.

The data from PubChem AID 1851 are available for download from the NIH through the PubChem website. I downloaded two files comprising the entire data set. The first, a structures file, contained the structural information encoded in the simplified molecular-input line-entry system (SMILES) format for each tested compound with corresponding Structure ID (SID) and Compound ID (CID) as assigned by NCBI. The second, also organized by SID and CID, contained all luminescent responses from the high-throughput screen and the fitted parameters summarized by the Activity Score. 
 

\section{Data Set Preparation and Molecular Descriptor Generation }

Both data files from PubChem AID 1851 were downloaded as comma separated value (.csv) files and merged on the SID column using Python's pandas library, a library for handling, manipulating and reshaping structured data.

Preprocessing typically begins with data cleaning, because raw data often contain anomalies, errors, or inconsistencies such as missing data, incomplete data, and invalid character values which may cause trouble in data analysis if left untreated. It is more complicated when data are collated from many formats that require harmonization and redundancy elimination~\cite{Nantasenamat2009}. Here, the source files needed little preprocessing.


\begin{figure}[ht]
  \centering
   \includegraphics[width=1\textwidth]{../img/Dataset_prep.png}
  \caption{Data Download and Preparation}
\end{figure}


\subsubsection{Molecular Descriptor Generation}
The merged file was loaded into a database in the Molecular Operating Environment (MOE). MOE's wash function removed salts and generated an energy-minimized conformation of each compound. Calculating the full suite of MOE 2-D descriptors added 186 columns of nominal, ordinal and continuous values to the database. The resulting master file was saved in .csv format.

The dataset was then split by isozyme into five separate files that each contained only the SID, the Activity Score, and the 186 MOE 2-D descriptors using a script in the Python programming language.

For each isozyme, non-inhibitors far outnumbered inhibitors. Because some later analyses require balanced classes, a Python script randomly shuffled each isozyme's non-inhibitors and kept only as many as there were inhibitors. 

\begin{figure}[H]
  \centering
   \includegraphics[width=.6\textwidth]{../img/Isozyme_Data_Prep.png}
  \caption{Isozyme Data Set Preparation}
\end{figure}

Next the script randomly shuffled the balanced datasets and split them into a training set and a test set. The training set held 80\% of compounds and the test set 20\%. The split was not stratified by class; the class ratio in each split was inspected to confirm it remained balanced.

For each use of randomness in computation, a seed number was set for the random number generator to ensure reproducible results.  

The balanced and split data sets are archived on Figshare for permanent, free and open access (http://dx.doi.org/10.6084/m9.figshare.1181846) and \\
(http://dx.doi.org/10.6084/m9.figshare.1066108). All subsequent analyses use these same splits for comparability.

\section{Data Preprocessing}

Data sets often contain redundant or noisy variables that make it harder for learning algorithms to discern meaningful patterns. Such multicollinearity can be exploited, or reduced to cut the computational cost of model construction~\cite{Nantasenamat2009}.

The variables in this data set varied widely in range and distribution, which causes problems for algorithms that use distance measures. Min-max normalization and z-score standardization address this. 

In min-max normalization, the minimum and maximum value of each variable is adjusted to a uniform range between 0 and 1 according to the following equation: $ X_{norm} = \frac{X - X_{min}}{X_{max} - X_{min}}  $ In z-score standardization, each variable is centered to mean 0 and scaled to unit variance: $ z = \frac{x - \mu}{\sigma}  $ \cite{Nantasenamat2009}

When data do not have a Gaussian (normal) distribution, simple transformations can make them more normal or symmetric. A common approach is a logarithmic transformation. This is typically performed on dependent variables such as the modeled biological/chemical properties of interest whereby IC50 may be transformed to logIC50 or -logIC50.~\cite{Nantasenamat2009}

\section{Modeling - Mapping Descriptors to Activity Data}

At this point, each isozyme had a complete data set, balanced between inhibitors and non-inhibitors and split into a training set and a test set. The preceding sections described data balancing and normalization; the following sections outline classification algorithms and validation procedures used in this study.

\begin{figure}[H]
  \centering
   \includegraphics[width=1\textwidth]{../img/Model_validation.png}
  \caption{Supervised Classification Model Training and Validation}
\end{figure}

\subsection{Binary QSAR in the Molecular Operating Environment}

The Binary QSAR approach described by Labute \cite{Labute1999} and as implemented in the 2011 version of MOE, classifies activity with a Bayesian model after reducing the descriptors to principal components. Principal component analysis (PCA) projects the data onto a smaller set of axes that point in the directions of greatest variance. The first principal component explains the most variance, the second the next most, and so on. Models are subsequently built with the original data projected onto any or all of these principal components.

Earlier work in Dr. Zheng's lab found that partial least squares (PLS) regression, a quantitative method in MOE that also begins with PCA, predicted poorly on this data set.

The training data were loaded into MOE and Binary QSAR was run from its menu-driven interface. The Activity Score threshold was set to 39, so compounds scoring 40 or above were labelled inhibitors and those scoring 39 or below non-inhibitors. The smoothing parameter was left at the default value of 0.25. MOE automatically performs principal component analysis on high-dimensional datasets before Binary QSAR; it offers no way to use the full suite of molecular descriptors untransformed.

For each isozyme, several models were built, each with a different number of principal components. Principal components are mutually orthogonal, so each captures a separate portion of the total variance. The assumption is that inclusion of more principal components leads to more of the variance being accounted for in a classification decision. However, since each principal component is a linear combination of all the variables, the benefits of dimensionality reduction come at the cost of interpretability. One aim of this study is to find the fewest principal components that still give good predictions.  For exploratory purposes, models with 2, 5, 10, 15, 20, 30, and 44 principal components were constructed. 

MOE models were written to .fit files and the model report saved as a .txt file.

\subsubsection{Modeling in Python}
Python is a dynamically typed, object-oriented, interpreted language. Its main strengths are rapid prototyping and a mature ecosystem of libraries, such as scikit-learn for machine learning. It is easy to learn, with abundant free learning materials online. The following classifiers were implemented with scikit-learn.

\subsection{\textit{k}-Nearest Neighbors (\textit{k}-NN)}
The \textit{k}-NN algorithm predicts the class of an object from the classes of its \textit{k} most similar training set objects~\cite{Lapins2013}. It finds the \textit{k} nearest points under a distance measure such as Euclidean distance and assigns the majority class among them.

\subsection{Random Forest (RF)}
A random forest is an ensemble of decision trees. Decision trees are made of nodes and branches. At each node, the data set is split based on the value of some attribute that is selected so that the instances of different classes are predominately moved to different branches. Classification passes each instance from the root node along the branches to a leaf node. To introduce diversity between the trees of a random forest, a small random subset of features is considered at each node of each tree. The forest classifies by majority vote of all trees~\cite{Lapins2013}.

\subsection{Support Vector Machines (SVM)}
An SVM finds the hyperplane that separates the classes with the greatest margin. SVMs use a kernel function to project the data into a higher-dimensional feature space before finding the hyperplane~\cite{Lapins2013}. The margin is the minimum distance from the data points to the hyperplane. The points closest to the hyperplane, called support vectors, determine its position.

\section{Model Evaluation and Validation}

The building blocks of a successful QSAR model are the accuracy of the input data, selection of appropriate descriptors and statistical tools, and most importantly validation of the developed model. Validation is the process where reliability and relevance of a procedure are established for a predetermined purpose. For QSAR models validation should target robustness, prediction performance and applicability domain of the models~\cite{Lapins2013}.

Statistical evaluation in QSAR modeling is essential to validate the model as well as to evaluate its predictive performance. Predictive performance can be assessed by dividing the data into a training set, used to construct the model, and a test set, used to evaluate it. Internal performance is typically assessed from the predictive performance within a training set, while external performance of a classifier can be assessed from the predictive performance against an independent test set that has never been seen by the training model. 

This study used two kinds of metric: overall accuracy, and the true positive and true negative rates. A common approach to model assessment in the exploratory phase is N-fold cross-validation: the data set is partitioned into N folds, and each fold in turn is left out of model construction and used for validation. For example, 5-fold cross-validation trains each model on four folds and validates it on the fifth. When samples are limited, leave-one-out cross-validation is the preferred approach. In that case, the number of folds is equal to the number of samples present in the data set~\cite{Nantasenamat2009}. Using these methods with the original training set allows the original test set to be preserved for final validation.

In this study, all Python models were assessed by 5-fold cross-validation on the training set. Each model was then fit on the full training set and used to predict the test set, which played no part in model fitting or feature scaling. The MOE models were built on the same training set and evaluated on the same test set, but without cross-validation.

\subsubsection{Confusion Matrix}

A confusion matrix (also called an error matrix) summarizes a classifier's performance. Each data point is counted by its actual class along one axis and by its predicted class along the other. For binary classification, this gives a $2 \times 2$ matrix.

\begin{tabular}{c @{\bfseries}r @{\hspace{0.7em}}c @{\hspace{0.4em}}c @{\hspace{0.7em}}l}
  \multirow{10}{*}{\parbox{1.1cm}{\bfseries\raggedleft Actual\\ Class}} & 
    & \multicolumn{2}{c}{\bfseries Prediction Class} & \\
  & & \bfseries p & \bfseries n &                     \\
  & p$'$ & \MyBox{True}{Positive} & \MyBox{False}{Negative} & \\[2.4em]
  & n$'$ & \MyBox{False}{Positive} & \MyBox{True}{Negative} & \\
  &  &  &  &
\end{tabular}

In this study, actual inhibitors (according to the assay) that the model predicts as inhibitors are true positives (TP), and actual non-inhibitors predicted as non-inhibitors are true negatives (TN). Actual inhibitors predicted as non-inhibitors are false negatives (FN), or Type II errors. Non-inhibitors predicted as inhibitors are false positives (FP), or Type I errors.

\subsubsection{Accuracy}

Accuracy is the fraction of instances classified correctly:
$$ ACC =\frac{(TP + TN)}{(TP + FP + TN + FN)} $$
where TP is the number of true positives, TN is the number of true negatives, FP is the number of false positives or over-predictions, and FN is the number of false negatives or missed predictions.

Accuracy is not an optimal measure of model performance if the data set is unbalanced (i.e. sizes of the classes are unequal) or if some errors are more costly than others (e.g. false negatives compared to false positives)~\cite{Lapins2013}. For this reason, all data sets were balanced to contain roughly equal numbers of inhibitors and non-inhibitors.

I also calculated the true positive rate (TPR) and true negative rate (TNR) for all models. TPR is the fraction of actual inhibitors predicted correctly:
$$ \mathrm{TPR} = \frac{TP}{TP + FN} $$
TPR is also called sensitivity; MOE calls it Accuracy on Actives.

TNR is the fraction of actual non-inhibitors predicted correctly:
$$ \mathrm{TNR} = \frac{TN}{TN + FP} $$
TNR is also called specificity; MOE calls it Accuracy on Inactives.

\subsubsection{Binary QSAR Test Set Evaluation}
The test sets were loaded and the washed structures were appended to the .csv files. All models were evaluated using MOE's menu-driven workflow and the classification probabilities were appended to the database file and saved as a .csv.  

In Microsoft Excel, predicted probabilities $\geq 0.5$ were classified as inhibitors and those $< 0.5$ as non-inhibitors. Confusion matrices, accuracy, TPR and TNR were then tabulated in the spreadsheet.

\subsubsection{Python Classifiers Using the Full Set of 2-D Descriptors}
The Python models omit PCA and use all 186 descriptors. For each isozyme, the Activity Score was converted to a binary response (1 for scores $\geq 40$, 0 otherwise), and all 186 descriptors served as predictors. Each classifier was combined with a standardization step (mean $0$, standard deviation $1$) in a scikit-learn pipeline, so the scaling parameters were estimated only from the data each model was fit on. Hyperparameters matched the scikit-learn 0.15 defaults: \textit{k}-NN with $k = 5$; a random forest of 10 trees considering $\sqrt{186}$ features at each split; and an SVM with a radial basis function kernel, $C = 1$ and $\gamma = 1/186$. Each model was first scored by 5-fold cross-validation on the training set, then fit on the full training set and used to predict the held-out test set. Accuracy, TPR and TNR were calculated from the resulting confusion matrices.

One notebook per isozyme runs the three classifiers through a shared evaluation module. Code and results are documented in Jupyter notebooks hosted on GitHub (https://github.com/5x5x5x5/CYP/), so the analysis is easy to verify and extend. Because scikit-learn gives every classifier the same interface, trying another algorithm requires only adding it to the list of classifiers.


\begin{figure}[!htbp]
  \centering
  \includegraphics[width=1\textwidth]{../img/supervised_learning_flowchart.png}
  \caption{Supervised Learning Workflow for Classification}
\end{figure}
